\chapter{The group theory of electron bands}\label{ch:bands}

\begin{watch}{\lec{7}}
Generators, graphically \ts{7}{196}{3:16} $\cdot$
tight binding for graphene \ts{7}{580}{9:40} $\cdot$
the point group and $\sh$ parity: $\upsigma$ and $\uppi$ bands \ts{7}{1083}{18:03} $\cdot$
little groups \ts{7}{1476}{24:36} $\cdot$
high-symmetry points and lines; why $K$ has $\Cv{3}$ \ts{7}{1728}{28:48} $\cdot$
the site-symmetry group \ts{7}{2104}{35:04} $\cdot$
orbital mixing and distorted orbitals \ts{7}{2421}{40:21} $\cdot$
the character formula for band states \ts{7}{3174}{52:54} $\cdot$
$\Gamma$ point \ts{7}{3551}{59:11} $\cdot$
$K$ point and its phases \ts{7}{3867}{64:27} $\cdot$
compatibility relations \ts{7}{4381}{73:01} $\cdot$
what symmetry cannot decide \ts{7}{5191}{86:31}
\end{watch}

This chapter applies the machinery to one concrete crystal. We label every $\uppi$ band
of graphene by small irreps, prove that the Dirac point is enforced by symmetry, and
introduce compatibility relations. They are the bridge to topology in
\cref{ch:topology}.

\section{Generators}\label{sec:generators}

A set of elements \emph{generates} $G$ if every element is a product of them and their
inverses. $C_n$ is generated by $C_n^+$; $\Cv{n}$ by $C_n^+$ and one mirror; $\Td$ by
$S_4$ and $C_3$ (or $S_4$ and a mirror that does not contain the $S_4$ axis). To check
that a set of functions transforms correctly under a representation, or that a matrix
commutes with a representation, it is enough to check the generators.

The lecture opens with a graphical check: mark a generic point on a stereographic
projection, apply the candidate generators repeatedly, and count the distinct images.
If their number equals $\abs{G}$, the set generates the group, and each element can be
read off as a word in the generators.

\section{A tight-binding model for graphene}

Carbon has the configuration $1s^2\,2s^2\,2p^2$. The $1s$ electrons are inert, so a
minimal tight-binding model uses the four valence orbitals $2s,2p_x,2p_y,2p_z$ on each
of the two atoms of the cell: $8\times8$ matrices $\mat{H}(\vec{k})$ and eight bands.
With nearest- and next-nearest-neighbour hopping such a model already reproduces
first-principles bands well near the Fermi level. Adding empty $3s$ orbitals improves
the upper bands.

\paragraph{The point group.} With the origin at a hexagon centre (\cref{fig:honeycomb}),
the honeycomb is invariant under rotations by multiples of $60^\circ$ and under six
vertical mirrors. They fall into two classes (\cref{fig:mirrors}):
\begin{itemize}
\item $\sv$: the three mirrors that \emph{contain} C--C bonds (lines at $30^\circ$,
      $90^\circ$, $150^\circ$ with our axes); they map each sublattice onto itself;
\item $\sd$: the three mirrors that \emph{bisect} C--C bonds (lines at $0^\circ$,
      $60^\circ$, $120^\circ$); they exchange sublattices A and B.
\end{itemize}
The in-plane point group is $\Cv{6}$ (plane group $p6mm$). Graphene is also invariant
under the reflection $\sh$ in its own plane, which doubles the group to $\Dh{6}$.
Since $\sh$ commutes with everything, we can use it once and for all:
$2s,2p_x,2p_y$ are even under $\sh$ and $2p_z$ is odd. The Hamiltonian cannot couple
states of opposite $\sh$ parity, so the $8\times8$ problem splits into the six
$\upsigma$ bands (bonding) and the two $\uppi$ bands from the $p_z$ orbitals, which contain
the Fermi level and the Dirac point. From now on we study the $\uppi$ bands with the group
$\Cv{6}$.

\begin{equation}\label{eq:C6vtable}
\begin{array}{c|cccccc}
\Cv{6} & E & C_2 & 2C_3 & 2C_6 & 3\sv & 3\sd\\\hline
\irr{A}_1 & 1&1&1&1&1&1\\
\irr{A}_2 & 1&1&1&1&-1&-1\\
\irr{B}_1 & 1&-1&1&-1&1&-1\\
\irr{B}_2 & 1&-1&1&-1&-1&1\\
\irr{E}_1 & 2&-2&-1&1&0&0\\
\irr{E}_2 & 2&2&-1&-1&0&0
\end{array}
\end{equation}
Which of the two $\irr{B}$ irreps is called $\irr{B}_1$ depends on which mirrors are
called $\sv$; with our choice $\irr{B}_1$ is even under the bond-containing mirrors.

\section{Little groups of graphene}

\begin{figure}[ht]
\centering
\begin{tikzpicture}[scale=1.25,>=Stealth]
  \draw[thick] (0:2)--(60:2)--(120:2)--(180:2)--(240:2)--(300:2)--cycle;
  \foreach \a in {30,90,150}{\draw[blue!60!black,dashed] (\a:-2.4)--(\a:2.4);}
  \foreach \a in {0,60,120}{\draw[red!60!black,dotted,thick] (\a:-2.4)--(\a:2.4);}
  \node[blue!60!black,font=\scriptsize] at (90:2.6) {$\sv$};
  \node[red!60!black,font=\scriptsize] at (60:2.65) {$\sd$};
  \fill (0,0) circle (0.05) node[above left,font=\small]{$\Gamma$};
  \fill (0:2) circle (0.05) node[right,font=\small]{$K$};
  \fill (-30:1.732) circle (0.05) node[below right,font=\small]{$M$};
  \draw[very thick,orange!80!black] (0,0)--(0:2) node[midway,above,font=\small]{$T$};
  \draw[very thick,orange!80!black] (0,0)--(-30:1.732) node[midway,below left,font=\small]{$\Sigma$};
  \draw[very thick,orange!80!black] (0:2)--(-30:1.732) node[midway,right,font=\small]{$T'$};
\end{tikzpicture}
\caption{The Brillouin zone of graphene with its high-symmetry points and lines. Blue
dashed: mirrors $\sv$ (containing C--C bonds in real space). Red dotted: mirrors $\sd$
(bisecting bonds). With $\vec{a}_1=a(1,0)$, the point $K=\tfrac23\vec{b}_1+\tfrac13\vec{b}_2
=(\tfrac{4\cpi}{3a},0)$ and $M=\tfrac12\vec{b}_1$.}
\label{fig:mirrors}
\end{figure}

The little co-group of each point or line is the set of point operations that map it
to an equivalent point:
\begin{equation}\label{eq:littlegroups}
\begin{array}{llll}
\toprule
\vec{k} & \bar{G}_{\vec{k}} & \text{elements} & \\
\midrule
\Gamma & \Cv{6} & \text{all 12} &\\
T\ (\Gamma K) & \Cs & E,\ \sd(0^\circ) &\\
K & \Cv{3} & E,\ C_3^\pm,\ 3\sd &\\
T'\ (KM) & \Cs & E,\ \sd(60^\circ) &\\
M & \Cv{2} & E,\ C_2,\ \sv(150^\circ),\ \sd(60^\circ) &\\
\Sigma\ (M\Gamma) & \Cs & E,\ \sv(150^\circ) &\\
\bottomrule
\end{array}
\end{equation}
Two remarks.
\begin{itemize}
\item At $K$ the rotation $C_3$ carries $K$ to another corner of the hexagon. The two
      corners differ by a reciprocal-lattice vector, so they label the \emph{same} irrep
      of $\Tgroup$ and $C_3\in\bar{G}_K$. For the same reason the six corners are only
      two inequivalent points, $K$ and $K'=-K$. This is why graphene has two Dirac
      points, not six.
\item The little co-group of a line is a subgroup of the little co-groups of its end
      points. This is what makes compatibility relations possible
      (\cref{sec:compat}).
\end{itemize}

\section{The site-symmetry group}\label{sec:sitesym}

\begin{definition}
The \emph{site-symmetry group} $G_{\vec{q}}$ of a point $\vec{q}$ is the set of
space-group elements that leave $\vec{q}$ fixed: $g\vec{q}=\vec{q}$ exactly (not up to a
lattice vector). It is always isomorphic to a point group. Points related by the space
group form an \emph{orbit}. Orbits whose site-symmetry groups are conjugate form a
\emph{Wyckoff position}.
\end{definition}
The crystal as a whole is invariant under $\Cv{6}$ about a hexagon centre, but an atom
does not see six-fold symmetry: along one direction it has a nearest neighbour, along
the opposite direction a third-nearest one. The site-symmetry group of a carbon atom
is $\Cv{3}$, generated by the three-fold rotation about the atom and the three $\sv$
mirrors through it. For $p6mm$ the Wyckoff positions with non-trivial site symmetry
are
\[
1a\ \text{(hexagon centres, }\Cv{6}),\qquad
2b\ \text{(the atoms, }\Cv{3}),\qquad
3c\ \text{(bond centres, }\Cv{2}).
\]

\paragraph{Atomic orbitals become irreps of the site group.} In the crystal the atomic
orbitals are no longer those of a free atom; they transform under irreps of
$G_{\vec{q}}$, obtained by subduction from $\Og{3}$ (\cref{sec:crystalfield}). At a
$\Cv{3}$ site, $s\to\irr{a}_1$ and $p\to\irr{a}_1\oplus\irr{e}$ with $p_z\in\irr{a}_1$
and $(p_x,p_y)\in\irr{e}$. The ``$p_z$ orbital'' of a tight-binding model of graphene
is really an $\irr{a}_1$ orbital of $\Cv{3}$: a $p_z$ orbital distorted by the crystal
field and mixed with other $\irr{a}_1$ orbitals.

This has practical consequences. A model built as if the orbitals were spherically
symmetric can miss couplings that are forbidden for free-atom orbitals but allowed by
the true, lower site symmetry. For example, the hoppings to two neighbours at equal
distance but in inequivalent directions need not be equal. Such an over-restricted
model may display spurious degeneracies or miss phases. Knowing the site-symmetry
irreps guarantees that the model is the most general one allowed by symmetry.

\section{Characters of band states}

Place one $\irr{a}_1$ orbital at each atom and form the two Bloch sums
$\Phi_{\mathrm{A}\vec{k}}$, $\Phi_{\mathrm{B}\vec{k}}$ of \eqref{eq:blochsum}. At each
$\vec{k}$ they span a two-dimensional representation of the little group, and
$\mat{H}(\vec{k})$ is a $2\times2$ matrix commuting with it. To find its small-irrep
content we need its character.

Let $g=\seitz{R}{\vec{0}}$ be in the little group of $\vec{k}$. If $g$ maps the atom at
$\vec{q}_\alpha$ onto an atom of the same sublattice, $g\vec{q}_\alpha=\vec{q}_\alpha+\vec{t}_\alpha$
with $\vec{t}_\alpha\in\Tgroup$, the Bloch sum is mapped onto itself up to a phase.
Otherwise it is mapped onto the other sublattice and does not contribute to the trace.
A short computation with \eqref{eq:blochsum} gives

\begin{keyresult}[Character of a band representation at $\vec{k}$]
\begin{equation}\label{eq:bandchar}
\chi_{\vec{k}}(g)=\sum_{\alpha\,:\,g\vec{q}_\alpha\equiv\vec{q}_\alpha}
\ee^{-\ii\vec{k}\cdot\vec{t}_\alpha}\;\chi_\rho(g_\alpha),
\qquad \vec{t}_\alpha=g\vec{q}_\alpha-\vec{q}_\alpha,
\end{equation}
where the sum runs over the atoms of the cell left invariant up to a lattice
translation, $g_\alpha=\seitz{E}{-\vec{t}_\alpha}g\in G_{\vec{q}_\alpha}$ is the
corresponding site-symmetry operation, and $\chi_\rho$ is the character of the orbital
irrep $\rho$ of the site group.
\end{keyresult}
This is the crystal version of $\chiM=\Ninv\chiV$ (\cref{sec:mechanical}): count invariant
atoms, weight each by the orbital character, and add the Bloch phase. For our
$\irr{a}_1$ orbitals $\chi_\rho=1$ and only the phases remain.

\subsection{The \texorpdfstring{$\Gamma$}{Gamma} point}
At $\vec{k}=0$ all phases are $1$ and $\chi_\Gamma(g)$ is the number of sublattices
mapped onto themselves: $2$ for $E$, $C_3$, $\sv$; $0$ for $C_2$, $C_6$, $\sd$ (these
exchange A and B). From \eqref{eq:C6vtable},
\[
\chi_\Gamma=(2,0,2,0,2,0)\qquad\Longrightarrow\qquad
\Gamma=\irr{A}_1\oplus\irr{B}_1
\]
(in the Bilbao labels used in the lecture, $\Gamma_1\oplus\Gamma_3$). Two
non-degenerate levels: the bonding combination $\Phi_{\mathrm{A}}+\Phi_{\mathrm{B}}$
($\irr{A}_1$) and the antibonding one $\Phi_{\mathrm{A}}-\Phi_{\mathrm{B}}$ ($\irr{B}_1$,
odd under the operations that exchange sublattices).

\subsection{The \texorpdfstring{$K$}{K} point}\label{sec:Kpoint}
Now phases matter. Only $E$ and $C_3^\pm$ keep the sublattices; the three $\sd$ in
$\bar{G}_K$ exchange them. Take $C_3^+$ about the origin. It maps
$\vec{q}_{\mathrm{A}}=\tfrac13(\vec{a}_1+\vec{a}_2)$ to $\vec{q}_{\mathrm{A}}-\vec{a}_1$ and
$\vec{q}_{\mathrm{B}}=\tfrac23(\vec{a}_1+\vec{a}_2)$ to $\vec{q}_{\mathrm{B}}-2\vec{a}_1$. With
$K=\tfrac23\vec{b}_1+\tfrac13\vec{b}_2$ we have $K\cdot\vec{a}_1=\tfrac{4\cpi}{3}$, so the
phases are
\[
\ee^{\ii K\cdot\vec{a}_1}=\ee^{4\cpi\ii/3}=\omega^{*},\qquad
\ee^{2\ii K\cdot\vec{a}_1}=\ee^{8\cpi\ii/3}=\omega,\qquad
\omega=\ee^{2\cpi\ii/3},
\]
and $\chi_K(C_3)=\omega+\omega^{*}=-1$. Physically: under the rotation, the Bloch sum
on sublattice A is multiplied by $\omega^{*}$, the one on B by $\omega$. Thus
\[
\chi_K=(2,-1,0)\ \text{on}\ (E,2C_3,3\sd)\qquad\Longrightarrow\qquad K=\irr{E}\quad(K_3).
\]
\begin{keyresult}[The Dirac point is enforced by symmetry]
At $K$ the two $\uppi$ states form the two-dimensional irrep $\irr{E}$ of $\Cv{3}$, so they
are degenerate for \emph{every} Hamiltonian with the symmetry of graphene. Whether the
dispersion near the degeneracy is linear is a further question, answered in
\cref{sec:kp}: it is.
\end{keyresult}

\subsection{The \texorpdfstring{$M$}{M} point}
$\bar{G}_M=\Cv{2}=\{E,C_2,\sv,\sd\}$. The $C_2$ and $\sd$ exchange sublattices, and
$\sv(150^\circ)$ keeps them with trivial phases, so $\chi_M=(2,0,2,0)$ and
$M=\irr{A}_1\oplus\irr{B}_1$, where in $\Cv{2}$ we call $\irr{B}_1$ the irrep that is
odd under $C_2$ and even under $\sv$.

\section{\texorpdfstring{$\vec{k}\cdot\vec{p}$}{k.p} Hamiltonians and the method of invariants\beyondmark}\label{sec:kp}
\sectionmark{The method of invariants}

Group theory tells us which bands must be degenerate at a high-symmetry point. The
\emph{method of invariants} tells us how they disperse nearby: it constructs the most
general effective Hamiltonian allowed by symmetry, order by order in the distance from
the point. It is the same idea as the invariant polarizability of \cref{sec:invariants}
and the invariant dynamical matrix of \cref{sec:wigner}.

\subsection{The effective Hamiltonian}
Let $\vec{k}_0$ be a high-symmetry point and let $\{\ket{n}\}$, $n=1,\dots,d$, be the states
of the (possibly degenerate) band multiplet of interest at $\vec{k}_0$, carrying a
representation $\mat{D}(g)$ of the little co-group $\bar{G}_{\vec{k}_0}$ (usually one small
irrep). For $\vec{k}=\vec{k}_0+\vec{q}$, projecting the Hamiltonian onto these states
(L\"owdin partitioning, which includes the other bands perturbatively) gives a
$d\times d$ matrix $\mat{H}(\vec{q})$, expanded in powers of $\vec{q}$:
\[
\mat{H}(\vec{q})=\mat{H}_0+\sum_iq_i\mat{H}_i+\sum_{ij}q_iq_j\mat{H}_{ij}+\dots
\]
At first order $\mat{H}_i\propto\bra{n}\hat{p}_i\ket{n'}$, hence the name
$\vec{k}\cdot\vec{p}$.

\subsection{The symmetry constraint}
For every $g$ in the little co-group, with rotational part $R_g$,
\begin{equation}\label{eq:kpconstraint}
\mat{D}(g)\,\mat{H}(R_g^{-1}\vec{q})\,\mat{D}(g)^{-1}=\mat{H}(\vec{q}).
\end{equation}
Expand $\mat{H}(\vec{q})=\sum_ac_a\,\mat{X}_a\,f_a(\vec{q})$, where the hermitian matrices
$\mat{X}_a$ span the $d^2$-dimensional space of $d\times d$ hermitian matrices and
$f_a$ are polynomials in $\vec{q}$. Both sets carry representations of the little
co-group: the matrices transform as $\mat{D}\otimes\mat{D}^{*}$ (by conjugation), the
polynomials of degree $n$ as the symmetric power $[\GV^n]$ restricted to the relevant
components of $\vec{q}$. Then \eqref{eq:kpconstraint} says that $\mat{H}$ is an
invariant, bilinear in matrices and polynomials. Hence
\begin{keyresult}[Method of invariants]
Decompose the matrices and the polynomials of each degree into irreps (with matched
bases). The most general $\mat{H}$ is a sum of one invariant, with a free real
coefficient, for each pair (matrix irrep, polynomial irrep) related as in
\cref{sec:countinv}. The number of free parameters at order $n$ equals the multiplicity
of $\irr{A}_1$ in $(\mat{D}\otimes\mat{D}^{*})\otimes[\GV^n]$.
\end{keyresult}
Additional antiunitary symmetries (time reversal, or time reversal combined with a
spatial operation that maps $\vec{k}_0$ to itself) add further constraints of the same
type.

\subsection{Example: the Dirac cone of graphene}
At $K$ the $\uppi$ states carry the two-dimensional irrep $\irr{E}$ of $\Cv{3}$
(\cref{sec:Kpoint}). Choose the band basis so that the matrices of $\irr{E}$ are those of
in-plane vectors, $\mat{D}(C_3^+)=$ rotation by $120^\circ$ and $\mat{D}(\sd)=\diag(1,-1)$
for the mirror that reflects $q_y\to-q_y$. The four hermitian $2\times2$ matrices
decompose as
\[
\mat{D}\otimes\mat{D}^{*}=\irr{E}\otimes\irr{E}=\irr{A}_1\oplus\irr{A}_2\oplus\irr{E}:
\qquad \one\ (\irr{A}_1),\quad \sigma_y\ (\irr{A}_2),\quad (\sigma_z,\sigma_x)\ (\irr{E}).
\]
Conjugation by a rotation by $\varphi$ rotates the pair $(\sigma_z,\sigma_x)$ by
$2\varphi$. For $\varphi=\pm120^\circ$ this is the same as rotating by $\mp120^\circ$, so
$(\sigma_z,-\sigma_x)$ transforms exactly like $(q_x,q_y)$. The polynomials decompose as:
degree $0$: $1$ ($\irr{A}_1$); degree $1$: $(q_x,q_y)$ ($\irr{E}$); degree $2$:
$q_x^2+q_y^2$ ($\irr{A}_1$) and $(q_x^2-q_y^2,\,2q_xq_y)$ ($\irr{E}$, rotating by $2\varphi$
like $(\sigma_z,\sigma_x)$). Pairing equal irreps gives, up to second order,
\begin{equation}\label{eq:dirac}
\mat{H}(\vec{q})=\varepsilon_0\one
+v\,(q_x\sigma_z-q_y\sigma_x)
+\alpha\,q^2\one
+\beta\bigl[(q_x^2-q_y^2)\sigma_z+2q_xq_y\sigma_x\bigr]+O(q^3).
\end{equation}
Consequences, all from symmetry:
\begin{itemize}
\item No constant $\sigma_y$ term: the mass term transforms as $\irr{A}_2$ and there is no
      $\irr{A}_2$ polynomial of degree $0$. The degeneracy at $K$ cannot be lifted without
      breaking the symmetry.
\item One linear invariant: the dispersion is linear and isotropic,
      $E-\varepsilon_0=\pm v\abs{\vec{q}}+O(q^2)$, a Dirac cone. The matrix
      $q_x\sigma_z-q_y\sigma_x$ is unitarily equivalent to the familiar
      $q_x\sigma_x+q_y\sigma_y$.
\item At second order, the $\beta$ term produces the three-fold \emph{trigonal warping}
      of the cone.
\item If the sublattices become inequivalent (as in hexagonal boron nitride), the
      relevant irrep at $K$ splits and a constant mass term, and a gap, become allowed.
\end{itemize}
The linear term exists because $\irr{E}\otimes\irr{E}$ contains the irrep of $\vec{q}$. When
that fails (for other groups or irreps), the crossing is quadratic or of higher order.
This is exactly the question asked in \lec{6}. Complete catalogues of such
$\vec{k}\cdot\vec{p}$ models for all space and magnetic groups have been built with this
method.

\paragraph{Further reading.}
G.\,L.\ Bir and G.\,E.\ Pikus, \emph{Symmetry and Strain-Induced Effects in Semiconductors}
(Wiley, 1974), the classic source of the method; R.\ Winkler, \emph{Spin--Orbit Coupling
Effects in Two-Dimensional Electron and Hole Systems} (Springer, 2003), with tables of
invariants; Dresselhaus, Dresselhaus and Jorio, the chapter on $\vec{k}\cdot\vec{p}$
perturbation theory.

\section{Compatibility relations}\label{sec:compat}

Along a line $\vec{k}$ varies continuously, and the little co-group is a subgroup of
those of the end points. A state at $\Gamma$ carrying an irrep $\Gamma^{(\mu)}$ evolves
continuously into the line, where it must carry the subduced representation
$\Gamma^{(\mu)}\!\downarrow\!\bar{G}_{\text{line}}$. These subductions are the
\emph{compatibility relations}. On every line $\bar{G}=\Cs$ with irreps $\irr{A}'$
(even under the mirror) and $\irr{A}''$ (odd). Reading the character of the
relevant mirror:
\[
\begin{array}{llll}
\toprule
\text{line (mirror)} & \text{from } \Gamma & \text{from } K & \text{from } M\\
\midrule
T\ (\sd) & \irr{A}_1\to\irr{A}',\ \irr{B}_1\to\irr{A}'' & \irr{E}\to\irr{A}'\oplus\irr{A}'' & \\
\Sigma\ (\sv) & \irr{A}_1\to\irr{A}',\ \irr{B}_1\to\irr{A}' & & \irr{A}_1\to\irr{A}',\ \irr{B}_1\to\irr{A}'\\
T'\ (\sd) & & \irr{E}\to\irr{A}'\oplus\irr{A}'' & \irr{A}_1\to\irr{A}',\ \irr{B}_1\to\irr{A}''\\
\bottomrule
\end{array}
\]
Each line can be reached from both ends, and the two subductions must agree. On $T$:
from $\Gamma$, $\irr{A}'\oplus\irr{A}''$; from $K$, $\irr{A}'\oplus\irr{A}''$. The same holds
on $\Sigma$ and $T'$. Such consistency checks catch almost every mistake. In practice
one never computes characters along lines; subduction from the end points gives them
for free.

\begin{figure}[ht]
\centering
\begin{tikzpicture}
\begin{axis}[width=11cm,height=7cm,xmin=0,xmax=9.9108,ymin=-4.2,ymax=3.2,
  xtick={0,4.1888,6.2832,9.9108},xticklabels={$\Gamma$,$K$,$M$,$\Gamma$},
  ylabel={$E/t$},xmajorgrids,grid style={gray!40},
  every axis plot/.append style={thick}]
\addplot[blue!70!black] table[x=x,y=lo]{graphene_pi.dat};
\addplot[red!70!black] table[x=x,y=hi]{graphene_pi.dat};
% labels at the high-symmetry points sit on the vertical grid lines: white background
\tikzset{pt/.style={font=\scriptsize,fill=white,inner sep=1.5pt}}
\node[pt,anchor=west]  at (axis cs:0.08,-3.95) {$\irr{A}_1$};
\node[pt,anchor=west]  at (axis cs:0.08,2.75) {$\irr{B}_1$};
\node[pt,anchor=south] at (axis cs:4.1888,0.6) {$\irr{E}$};
\node[pt,anchor=north] at (axis cs:6.2832,-1.0) {$\irr{B}_1$};
\node[pt,anchor=south] at (axis cs:6.2832,1.4) {$\irr{A}_1$};
% labels on the lines: each sits on the free side of its (monotonic) curve
\tikzset{ln/.style={font=\scriptsize,inner sep=2pt}}
\node[ln,blue!70!black,anchor=north west] at (axis cs:2.13,-2.06) {$\irr{A}'$};
\node[ln,red!70!black,anchor=south west]  at (axis cs:2.13,1.88)  {$\irr{A}''$};
\node[ln,blue!70!black,anchor=north east] at (axis cs:5.27,-0.51) {$\irr{A}''$};
\node[ln,red!70!black,anchor=south east]  at (axis cs:5.27,0.99)  {$\irr{A}'$};
\node[ln,blue!70!black,anchor=north east] at (axis cs:8.13,-2.48) {$\irr{A}'$};
\node[ln,red!70!black,anchor=south east]  at (axis cs:8.13,2.05)  {$\irr{A}'$};
\end{axis}
\end{tikzpicture}
\caption{The $\uppi$ bands of graphene in a tight-binding model with nearest- ($t$) and
next-nearest-neighbour ($t'=0.1\,t$) hopping, labelled by small irreps of $\Cv{6}$ and
its subgroups. The labels at $\Gamma$ and $M$ were obtained from the eigenvectors;
symmetry alone fixes only which pair of labels appears at each point.}
\label{fig:grapheneTB}
\end{figure}

\section{What symmetry decides and what it does not}

Symmetry fixes the \emph{set} of irreps at each point: $\{\irr{A}_1,\irr{B}_1\}$ at
$\Gamma$, $\{\irr{E}\}$ at $K$, $\{\irr{A}_1,\irr{B}_1\}$ at $M$. It does not fix their
energetic order. In \cref{fig:grapheneTB} the $\irr{A}_1$ state is below at $\Gamma$ but
above at $M$. The model produced this, not symmetry, and different hoppings could
reorder them. The compatibility relations then decide how the labels are connected.
The order of the irreps at high-symmetry points, combined with the compatibility
relations, is exactly the information that topological quantum chemistry uses.

Irrep labels are also what selection rules need: optical transitions between bands at
$\vec{k}$ are allowed or forbidden according to the small irreps involved
(\cref{sec:raman}).

\exercisesection

\begin{exercise}\label{ex:7:littlegroups}
Verify the little co-groups \eqref{eq:littlegroups}. In particular show that
$\sd(60^\circ)$ maps $K$ to an equivalent point and that $\sv(150^\circ)$ fixes $M$.
\end{exercise}

\begin{solution}
$\sd(60^\circ)$ reflects the direction $0^\circ$ of $K$ to $120^\circ$, i.e.\ to $C_3^+K$,
which is equivalent to $K$ (\cref{ex:6:recip}). $M$ lies at $-30^\circ$, on the mirror line
at $150^\circ\equiv-30^\circ$, so $\sv(150^\circ)$ fixes it exactly. By contrast
$\sv(30^\circ)$ sends $M$ to $90^\circ$, i.e.\ to $\vec{b}_2/2$, and
$\vec{b}_2/2-\vec{b}_1/2$ is not a reciprocal vector. The lines inherit the mirror that
contains them: $\sd(0^\circ)$ for $T$, $\sv(150^\circ)$ for $\Sigma$, and $\sd(60^\circ)$
for $T'$ (it maps the segment $KM$ onto an equivalent one).
\end{solution}

\begin{exercise}[the $\upsigma$ bands]\label{ex:7:sigma}
The $\upsigma$ bands come from $s$ ($\irr{a}_1$) and $(p_x,p_y)$ ($\irr{e}$) orbitals at the
$\Cv{3}$ sites. Use \eqref{eq:bandchar} to show that the six $\upsigma$ states decompose as
\[
\Gamma:\ \irr{A}_1\oplus\irr{B}_1\oplus\irr{E}_1\oplus\irr{E}_2,\qquad
K:\ \irr{A}_1\oplus\irr{A}_2\oplus2\,\irr{E},\qquad
M:\ 2\,\irr{A}_1\oplus\irr{A}_2\oplus2\,\irr{B}_1\oplus\irr{B}_2 .
\]
How many symmetry-enforced degeneracies are there at $\Gamma$ and at $K$?
\end{exercise}

\begin{solution}
Each atom carries $\irr{a}_1\oplus\irr{e}$ with site characters
$\chi_\rho(E)=3$, $\chi_\rho(C_3)=1-1=0$, $\chi_\rho(\sv)=1+0=1$.
$\Gamma$: sublattices kept by $E,C_3,\sv$ with trivial phases,
$\chi_\Gamma=(6,0,0,0,2,0)$ on $(E,C_2,2C_3,2C_6,3\sv,3\sd)$. The reduction gives
$m_{\irr{A}_1}=m_{\irr{B}_1}=\tfrac1{12}(6+6)=1$, $m_{\irr{E}_1}=m_{\irr{E}_2}=\tfrac1{12}\,12=1$,
others $0$.
$K$: $\chi_K(C_3)=(\omega+\omega^{*})\cdot0=0$, $\sd$ exchanges sublattices, so
$\chi_K=(6,0,0)$ and $K=\irr{A}_1\oplus\irr{A}_2\oplus2\,\irr{E}$.
$M$: $\chi_M=(6,0,2,0)$ on $(E,C_2,\sv,\sd)$, giving
$2\,\irr{A}_1\oplus\irr{A}_2\oplus2\,\irr{B}_1\oplus\irr{B}_2$.
Enforced degeneracies: two doubly degenerate levels at $\Gamma$ ($\irr{E}_1$, $\irr{E}_2$) and
two at $K$ ($2\,\irr{E}$); none at $M$.
\end{solution}

\begin{exercise}[nearest-neighbour model]\label{ex:7:NN}
With only nearest-neighbour hopping $t$, $\mat{H}(\vec{k})=-t\begin{pmatrix}0&f\\f^{*}&0\end{pmatrix}$
with $f(\vec{k})=1+\ee^{-\ii\vec{k}\cdot\vec{a}_1}+\ee^{-\ii\vec{k}\cdot\vec{a}_2}$ (Bloch sums
\eqref{eq:blochsum}). Compute the energies at $\Gamma$, $K$ and $M$ and identify the
irrep of each eigenvector by applying $C_2$.
\end{exercise}

\begin{solution}
$\Gamma$: $f=3$, $E=\mp3t$. The lower state $(1,1)/\sqrt2$ is invariant under all
operations: $\irr{A}_1$; the upper one $(1,-1)/\sqrt2$ is odd under the operations
exchanging sublattices: $\irr{B}_1$.
$K$: $K\cdot\vec{a}_1=\tfrac{4\cpi}{3}$, $K\cdot\vec{a}_2=\tfrac{2\cpi}{3}$, so
$f=1+\omega+\omega^{*}=0$: $E=0$, doubly degenerate ($\irr{E}$).
$M$: $M\cdot\vec{a}_1=\cpi$, $M\cdot\vec{a}_2=0$, $f=1-1+1=1$, $E=\mp t$. $C_2$ maps
$\vec{q}_{\mathrm{A}}\to-\vec{q}_{\mathrm{A}}=\vec{q}_{\mathrm{B}}-(\vec{a}_1+\vec{a}_2)$, so
$\hat{O}_{C_2}\Phi_{\mathrm{A}}=\ee^{\ii M\cdot(\vec{a}_1+\vec{a}_2)}\Phi_{\mathrm{B}}=-\Phi_{\mathrm{B}}$ and
likewise $\hat{O}_{C_2}\Phi_{\mathrm{B}}=-\Phi_{\mathrm{A}}$. The lower state $(1,1)$ has $C_2$
eigenvalue $-1$ ($\irr{B}_1$), the upper $(1,-1)$ has $+1$ ($\irr{A}_1$), as in
\cref{fig:grapheneTB}.
\end{solution}

\begin{exercise}\label{ex:7:compat}
Write the compatibility relations for the $\upsigma$ bands of \cref{ex:7:sigma} along
the line $T$, and check that the subductions from $\Gamma$ and from $K$ agree.
\end{exercise}

\begin{solution}
On $T$ the mirror is $\sd$. From $\Gamma$: $\irr{A}_1\to\irr{A}'$, $\irr{B}_1\to\irr{A}''$,
$\irr{E}_1\to\irr{A}'\oplus\irr{A}''$, $\irr{E}_2\to\irr{A}'\oplus\irr{A}''$: in total
$3\irr{A}'\oplus3\irr{A}''$. From $K$: $\irr{A}_1\to\irr{A}'$, $\irr{A}_2\to\irr{A}''$,
$2\,\irr{E}\to2(\irr{A}'\oplus\irr{A}'')$: again $3\irr{A}'\oplus3\irr{A}''$. The two subductions
agree.
\end{solution}
